\chapter*{Abstract}

Humans navigate complex environments by abstracting situational details into context states, enabling efficient perception and action. The Free Energy Principle (FEP) provides a framework for understanding this behavior, positing that agents minimize expected free energy, which naturally gives rise to both goal-directed action and information-seeking (curiosity). While epistemic behavior regarding immediate environmental states has been well-studied, contextual information-seeking—where agents actively resolve uncertainty about abstract situational rules—remains underexplored. To investigate these cognitive mechanisms, this thesis develops a computational behavioral model, the Context-GFE agent, capable of actively seeking information about an abstract context state. Grounded in Active Inference and formalized using Forney-style Factor Graphs (FFG), message passing algorithms are derived to enable the unified inference of past, present, and future belief distributions. The newly proposed Context-GFE agent is evaluated against three baseline models (BFE, Context-BFE, and GFE) in a simulated T-maze environment. Results demonstrate that the Context-GFE agent successfully resolves contextual uncertainty by planning ahead to receive informative observations, balancing epistemic and pragmatic drives. Overall, this model provides a mathematically grounded and testable hypothesis for studying contextual abstraction and cognitive control, facilitating future empirical validation through Bayesian model comparison.